\documentclass[12pt]{article}

\usepackage[a4paper,margin=1in]{geometry}
\usepackage{amsmath,amssymb}
\usepackage{setspace}
\usepackage{enumitem}

\setstretch{1.2}

\begin{document}

\begin{center}
\textbf{\Large CSE400 – Fundamentals of Probability in Computing} \\[4pt]
\textbf{Lecture 6: Discrete Random Variables, Expectation and Problem Solving} \\[8pt]
Instructor: Dhaval Patel, PhD \\
Date: January 22, 2025
\end{center}

\vspace{1em}

\noindent \textbf{L6\_S1\_A}

\section{Random Variables (RVs)}

\subsection{Motivation and Concept}

A random variable (RV) is a function that maps outcomes of a random experiment (sample space) to numerical values.

The distribution of a random variable can be visualized using a bar diagram.

The x-axis represents the values that the random variable can take.

The height of the bar at value $a_i$ represents:
\[
\Pr[X = a_i]
\]

Each probability is computed by finding the probability of the corresponding event in the sample space.

\section{Types of Random Variables}

\subsection{Discrete Random Variables}

A discrete random variable has the following properties:
\begin{itemize}
    \item Countable support
    \item Defined using a probability mass function (PMF)
    \item Probabilities are assigned to single values
    \item Each possible value has strictly positive probability
\end{itemize}

\subsection{Continuous Random Variables}

A continuous random variable has the following properties:
\begin{itemize}
    \item Uncountable support
    \item Defined using a probability density function (PDF)
    \item Probabilities are assigned to intervals of values
    \item Each specific value has zero probability
\end{itemize}

\section{Example: Random Variable from Coin Tossing}

\textbf{Example 1: Tossing 3 Fair Coins}

Experiment: Toss 3 fair coins.  
Let $Y$ denote the number of heads observed.

Possible values of $Y$:
\[
Y \in \{0,1,2,3\}
\]

\textbf{Probability Computation:}

\[
P(Y=0) = P(t,t,t) = \frac{1}{8}
\]

\[
P(Y=1) = P(t,t,h),(t,h,t),(h,t,t) = \frac{3}{8}
\]

\[
P(Y=2) = P(h,h,t),(h,t,h),(t,h,h) = \frac{3}{8}
\]

\[
P(Y=3) = P(h,h,h) = \frac{1}{8}
\]

\textbf{Verification of Total Probability:}

Since $Y$ must take one of the values $0$ through $3$,
\[
1 = P\left( \bigcup_{i=0}^{3} \{Y=i\} \right)
= \sum_{i=0}^{3} P(Y=i)
\]

\section{Probability Mass Function (PMF)}

\subsection{Definition}

A random variable that can take at most a countable number of values is called discrete.

Let $X$ be a discrete random variable with range:
\[
R_X = \{x_1, x_2, x_3, \dots\}
\]
where the range is finite or countably infinite.

The function
\[
P_X(x_k) = P(X = x_k), \quad k = 1,2,3,\dots
\]
is called the Probability Mass Function (PMF) of $X$.

\subsection{PMF Property}

Since $X$ must take one of the values $x_k$, we must have:
\[
\sum_{k=1}^{\infty} P_X(x_k) = 1
\]

\section{PMF Example: Given Functional Form}

\textbf{Example}

The PMF of a random variable $X$ is given by:
\[
p(i) = c \frac{\lambda^i}{i!}, \quad i = 0,1,2,\dots
\]
where $\lambda > 0$.

\textbf{Tasks:}
\begin{enumerate}
    \item Find $P(X=0)$
    \item Find $P(X>2)$
\end{enumerate}

\subsection*{Step 1: Determine Constant $c$}

Since probabilities must sum to 1:
\[
\sum_{i=0}^{\infty} p(i) = 1
\]

Substitute:
\[
c \sum_{i=0}^{\infty} \frac{\lambda^i}{i!} = 1
\]

Using the known series:
\[
e^{\lambda} = \sum_{i=0}^{\infty} \frac{\lambda^i}{i!}
\]

We get:
\[
c e^{\lambda} = 1 \Rightarrow c = e^{-\lambda}
\]

\subsection*{Step 2: Compute $P(X=0)$}

\[
P(X=0) = e^{-\lambda} \frac{\lambda^0}{0!} = e^{-\lambda}
\]

\subsection*{Step 3: Compute $P(X>2)$}

\[
P(X>2) = 1 - P(X \le 2)
\]

\[
= 1 - [P(X=0) + P(X=1) + P(X=2)]
\]

\[
= 1 - \left(e^{-\lambda} + \lambda e^{-\lambda} + \frac{\lambda^2}{2} e^{-\lambda}\right)
\]

\section{Bayes’ Theorem (Recap)}

\subsection{Conditional Probability Identity}

Using:
\[
\Pr(A|B) = \Pr(B|A)\Pr(A)
\]

We obtain:
\[
\Pr(B_i|A) =
\frac{\Pr(A|B_i)\Pr(B_i)}
{\sum_{j=1}^{n} \Pr(A|B_j)\Pr(B_j)}
\]

This is known as Bayes’ Formula (Proposition 3.1).

\subsection{Terminology}

$\Pr(B_i)$: A priori probability  
(probability formed from self-evident or presupposed models)

$\Pr(B_i|A)$: Posterior probability  
(probability computed after observing event $A$)

\section{Bayes’ Theorem – Example 1: Auditorium with 30 Rows}

\textbf{Problem Setup}

An auditorium has 30 rows.

Row 1 has 11 seats  
Row 2 has 12 seats  
Row 3 has 13 seats  
\[
\vdots
\]
Row 30 has 40 seats

A door prize is awarded by:
\begin{enumerate}
    \item Randomly selecting a row (each row equally likely)
    \item Randomly selecting a seat within that row
\end{enumerate}

\textbf{Task 1: Probability that Seat 15 is Selected Given Row 20}

\[
P(S_{15} | R_{20}) = \frac{1}{30}
\]

\textbf{Task 2: Probability that Row 20 was Selected Given Seat 15}

Using Bayes’ rule:
\[
P(R_{20} | S_{15}) =
\frac{P(S_{15} | R_{20})P(R_{20})}{P(S_{15})}
\]

Where:
\[
P(S_{15}) = \sum_{k=5}^{30} P(S_{15} | R_k) P(R_k)
\]

Numerical evaluation (as shown in lecture):
\[
P(S_{15}) = 0.0342
\]

Final result:
\[
P(R_{20} | S_{15}) =
\frac{\frac{1}{30} \cdot \frac{1}{30}}{0.0342}
= 0.0325
\]

\section{Bayes’ Theorem – Example 2: Communication System (Receiver RX)}

A communication system sends binary data (0 or 1).

The receiver (RX) may make mistakes:
\begin{itemize}
    \item A 0 may be detected as 1
    \item A 1 may be detected as 0
\end{itemize}

The system is described using a set of conditional probabilities, and Bayes’ theorem is used to compute posterior probabilities of transmitted symbols given received symbols.

\end{document}
